% TOP_PROBLEM: {"id": 2302068, "title": "Research Problems in Function Theory — Problem 2.68", "category_id": 8, "category": {"id": 8, "name": "analysis", "display_name": "Analysis", "description": "Limits, continuity, calculus, and function theory.", "slug": "analysis", "order_index": 8, "created_at": "2026-07-31T15:26:25.671Z"}, "status": "partially_solved", "problem_number": "AMR-022-2068", "set_id": 13}
% FIRST_POSTED: 2026-10-02
% ENTRY_KIND: statement_only
% UnsolvedMath ID 2302068; source catalogue code AMR-022-2068
% !TeX program = lualatex
% Definition and sources only; this file is not an LLM solution attempt.
\documentclass[11pt,a4paper]{article}
\usepackage[margin=25mm]{geometry}
\usepackage{fontspec}
\setmainfont{lmroman10-regular.otf}[BoldFont=lmroman10-bold.otf,
 ItalicFont=lmroman10-italic.otf,BoldItalicFont=lmroman10-bolditalic.otf]
\usepackage{amsmath,amssymb,mathtools}
\usepackage{unicode-math}
\setmathfont{latinmodern-math.otf}
\AtBeginDocument{\let\setminus\smallsetminus}
\usepackage{xurl,hyperref}
\hypersetup{colorlinks=true,urlcolor=blue}
\setcounter{secnumdepth}{0}
\setlength{\parindent}{0pt}
\setlength{\parskip}{5pt}
\setlength{\emergencystretch}{3em}
\begin{document}
\section*{Research Problems in Function Theory — Problem 2.68}\label{2302068}
{\small UnsolvedMath ID 2302068 · AMR-022-2068 · Analysis · AMR Open Problem Lists}
\subsection{Definitions and mathematical statement}
Let $\rho>0$ and let $f$ be an entire function with
\[
\log M(r,f)\le(1+o(1))r^\rho\quad(r\to\infty),
\qquad M(r,f)=\max_{|z|=r}|f(z)|.
\]
Suppose a continuous curve $\Gamma$ tends to infinity and, for some fixed $\alpha\in[-1,1)$, satisfies
$\log|f(z)|\le(\alpha+o(1))|z|^\rho$ as $z\to\infty$ along $\Gamma$.
For $\varepsilon>0$ define the angular measure, in radians,
\[
E(r,\varepsilon)=\bigl|\{\theta\in[0,2\pi):
\log|f(re^{i\theta})|\le(1-\varepsilon)r^\rho\}\bigr|.
\]
Zeros are included by the convention $\log0=-\infty$. The conjecture is
\[
\lim_{\varepsilon\downarrow0}\limsup_{r\to\infty}E(r,\varepsilon)
\ge \frac2\rho\arccos\alpha.
\]
The outer limit exists by monotonicity as $\varepsilon$ decreases. This specifies the joint limiting notation used in the source. The same curve supplies the small-value hypothesis at all sufficiently large points on it; a sequence of unrelated small-value points is insufficient. The proposed conclusion measures how much of some large circles must consequently lie below the near-maximal growth scale.
\subsection{Short English statement}
Does unusually small growth along one curve force a quantitatively large set of small-growth angles on arbitrarily large circles?
\subsection{Sources}
\begin{itemize}
\item[S1] ULAM.AI and the UnsolvedMath Contributors, supplied problems.json, record 2302068 (AMR-022-2068), AMR Open Problem Lists. Catalogue identity and source transcription; CC BY 4.0.
\url{https://huggingface.co/datasets/ulamai/UnsolvedMath}.
\item[S2] Hayman - Research Problems in Function Theory (2018), Problem 2.68. Original source indexed by AMR-022-2068.
\url{https://arxiv.org/abs/1809.07200}.
\end{itemize}
\end{document}
